% Wrapper-free contribution. Requires amsmath, amssymb, amsthm and theorem environments.
% All mathematical assertions below are proved in this file.
\section{Embedded arithmetic-site points and common primitive shell pairs}
\label{sec:connes-primitive}

\paragraph{Primary source and precise specialization.}
Alain Connes and Caterina Consani, \emph{Cyclic theory and the pericyclic
category}, in \emph{Cyclic Cohomology at 40: Achievements and Future Prospects},
Proceedings of Symposia in Pure Mathematics 105 (2023), pp.~103--122,
introduce the divisor category $\mathcal R$ in \S7.4.
The chapter DOI is \texttt{10.1090/pspum/105/01897}.
The category has one morphism $n\to m$ when $m\mid n$.
Their Lemma~7.7, printed p.~120, identifies its topos points with embedded
groups $\mathbb Z\subset H\subset\mathbb Q$; Proposition~7.8, printed
p.~121, calculates the further passage to an abstract ordered group.
The supplied volume's corresponding PDF pages are 131 and 132.
Both complete Connes--Consani chapters, printed pp.~83--102 and 103--122,
were read before making this selection.  The following is a finite
specialization with an explicit primitive-pair return, proved independently
here.  No claim that the source already contains the shell theorem is made.

For a positive integer $N=\prod_\ell\ell^{e_\ell}$ start with the actual
finite box
\[
 \mathcal B_N=\prod_{\ell\mid N}\{0,\ldots,e_\ell\},\qquad
 \mathcal D_N=\left\{\prod_{\ell\mid N}\ell^{j_\ell}:
             (j_\ell)\in\mathcal B_N\right\}.
\]
Thus $\mathcal D_N$ is precisely the positive divisors of $N$.
The functor $F_N:\mathcal R\to\mathrm{Sets}$ sends $n$ to one point
when $n\mid N$ and to the empty set otherwise.  It is well defined by
transitivity of divisibility.  Its category of elements is nonempty,
and the least common multiple of two divisors of $N$ is a divisor of $N$
mapping to both; there are no distinct parallel arrows.  These are the
filtering conditions used in the source's Lemma~7.7.  Its associated
embedded group is
\[
 H_N=\bigcup_{n\in\mathcal D_N}\frac1n\mathbb Z\subset\mathbb Q.
\]
This is a specified construction from the box, with an exact return map;
the box is not discarded in favour of an unmarked group.

\begin{proposition}[Finite embedded points and exact divisor return]
\label{prop:connes-finite-return}
One has
\[
 H_N=\frac1N\mathbb Z,\qquad
 \mathcal D_N=\{d\in\mathbb Z_{>0}:1/d\in H_N\}.
\]
For positive $N,M$, if $g=\gcd(N,M)$ and $L=\operatorname{lcm}(N,M)$,
then
\[
 H_N\cap H_M=H_g,\qquad H_N+H_M=H_L.
\]
In particular there is an exact sequence of abelian groups
\[
 0\longrightarrow H_g\xrightarrow{t\mapsto(t,-t)}H_N\oplus H_M
 \xrightarrow{(x,y)\mapsto x+y}H_L\longrightarrow0.
\]
Every fibre of the last arrow is a translate of the displayed image of
$H_g$.  Also $H_N\subset H_M$ holds exactly when $N\mid M$.
\end{proposition}
\begin{proof}
For $n\mid N$, $1/n=(N/n)/N$, and $n=N$ occurs in the union.
Thus $H_N=(1/N)\mathbb Z$.  The equality $1/d=k/N$ with $k\in\mathbb Z$
is equivalent to $dk=N$, and for $d>0$ it is equivalent to $d\mid N$.
This proves the exact return, including all prime valuations and their bounds.
An $x\in H_N\cap H_M$ has $Nx,Mx\in\mathbb Z$.
Choose integers $r,s$ with $rN+sM=g$; then $gx\in\mathbb Z$, proving
$x\in H_g$.  Conversely $g\mid N,M$ implies $H_g\subset H_N\cap H_M$.
The sum is generated by $1/N=(M/g)/L$ and $1/M=(N/g)/L$.
Since $\gcd(M/g,N/g)=1$, a Bezout combination gives $1/L$, proving the
sum formula.  The kernel and its fibres follow by subtracting two lifts
of the same sum: a difference in the kernel is $(t,-t)$ with
$t\in H_N\cap H_M$.  Finally $H_N\subset H_M$ implies $1/N\in H_M$,
so the exact divisor return gives $N\mid M$; the converse follows from
$1/N=(M/N)/M$.
\end{proof}


For such an embedded group define its primitive reciprocal pairs by
\[
 \mathcal P(H)=\{(1/m,1/n)\in H^2:
                  m,n\in\mathbb Z_{>0},\ \gcd(m,n)=1\}.
\]
In the original finite box this means that for each prime the two
valuations lie between zero and its exponent in $N$, and at least one
valuation is zero.  Thus the primitivity decoration retains an explicit
exclusive-use condition on each finite exponent budget.

Fix the same prime $p=12h+1$ throughout, and let
$a,b\in\{3h+1,\ldots,9h\}$.  Keep
\[
 S_a=pa,\quad R_a=4a-p,\qquad S_b=pb,\quad R_b=4b-p.
\]
For each shell $t\in\{a,b\}$ define the ordered marked pair set
\[
 \mathcal Q_t=\{(m,n)\in\mathbb Z_{>0}^2:
 m\mid S_t,\ n\mid S_t,\ \gcd(m,n)=1,\ R_t\mid m+n\}.
\]
Both $R_t$ are positive odd integers at least $3$.
Also $\gcd(R_t,S_t)=1$: a common prime divisor with $t$ would divide
$p$ but $0<t<p$; a common divisor with $p$ would divide $4t$, whereas
$p$ is an odd prime not dividing $t$.

\begin{theorem}[Embedded intersection with full shell marks]
\label{thm:connes-common-pairs}
Set
\[
 d=\gcd(a,b),\qquad L_R=\operatorname{lcm}(R_a,R_b).
\]
Then
\[
 H_{S_a}\cap H_{S_b}=H_{pd},
\]
and reciprocal embedding gives the exact bijection
\[
 \mathcal Q_a\cap\mathcal Q_b
 \longleftrightarrow
 \{(1/m,1/n)\in\mathcal P(H_{pd}):L_R\mid m+n\}.
\]
Its map is $(m,n)\mapsto(1/m,1/n)$ and its inverse is coordinatewise
reciprocal.  In integer coordinates the common pair set is precisely
\[
 \{(m,n):m\mid p\gcd(a,b),\ n\mid p\gcd(a,b),\
          \gcd(m,n)=1,\ L_R\mid m+n\}.
\]
The fibres of this correspondence are singletons.  For noncoprime shell
moduli the retained sum modulus is exactly
$L_R=R_aR_b/\gcd(R_a,R_b)$.
\end{theorem}
\begin{proof}
Since $p$ is the same factor in both $S_a,S_b$,
$\gcd(S_a,S_b)=p\gcd(a,b)=pd$.  Apply
Proposition~\ref{prop:connes-finite-return} to get the intersection.
An integer $m$ divides both $S_a,S_b$ exactly when $1/m$ lies in both
embedded groups, equivalently in $H_{pd}$; the same holds for $n$.
The primitivity condition is unchanged under this comparison because
the integer labels $m,n$ are recovered exactly by reciprocals.
For a positive integer $z$, divisibility by both $R_a$ and $R_b$ is
equivalent to divisibility by their least common multiple: for each prime
its valuation must be at least the maximum of the two exponents.
Applying this to $z=m+n$ proves the mark and the displayed formula for
$L_R$.  Both reciprocal maps undo each other, proving bijectivity and
the asserted fibres.
\end{proof}

\begin{proposition}[Explicit return of every common pair to both witnesses]
\label{prop:connes-pair-witness}
For $(m,n)$ in the common pair set and $t=a,b$, define
\[
 k_t=\frac{S_t}{mn}\frac{m+n}{R_t},\qquad
 B_t=k_tm,\quad C_t=k_tn.
\]
These are positive integers and
\[
 \frac4p=\frac1t+\frac1{B_t}+\frac1{C_t},\qquad
 \frac{B_t}{C_t}=\frac mn.
\]
Conversely every ordered positive integral denominator pair $(B_t,C_t)$
at shell $t$ determines exactly one pair in $\mathcal Q_t$ by
\[
 k_t=\gcd(B_t,C_t),\qquad m=B_t/k_t,\quad n=C_t/k_t.
\]
Thus an ordered rational ratio occurs at both shells exactly when its
unique primitive pair lies in the common set of
Theorem~\ref{thm:connes-common-pairs}.  The two shell witness pairs for
such a ratio are then exactly the displayed formulas.
\end{proposition}
\begin{proof}
The coprimality of $m,n$ and their separate divisibility into $S_t$ give
$mn\mid S_t$: at each prime at most one of the two positive valuations
occurs, and that valuation is at most the exponent in $S_t$.
Also $R_t\mid m+n$.  Hence $k_t,B_t,C_t$ are positive integers, and
\[
 \frac1{B_t}+\frac1{C_t}
 =\frac{m+n}{k_tmn}=\frac{R_t}{S_t}=\frac4p-\frac1t.
\]
The ratio equality is immediate from the positive factor $k_t$.
Conversely the reciprocal identity gives
$R_t B_t C_t=S_t(B_t+C_t)$.  Set $k=\gcd(B_t,C_t)$,
$m=B_t/k$, and $n=C_t/k$; the definitions give $\gcd(m,n)=1$ and
\[
 R_t k mn=S_t(m+n).
\]
No prime dividing $mn$ divides $m+n$, so $mn\mid S_t$.  Dividing by
$mn$ gives $R_tk=(S_t/(mn))(m+n)$; as
$\gcd(R_t,S_t/(mn))=1$, it follows that $R_t\mid m+n$.
Thus $(m,n)\in\mathcal Q_t$, and the same equation forces exactly
$k=(S_t/(mn))((m+n)/R_t)$.
For completeness, uniqueness of the primitive pair for a positive ratio
follows as well: if $m/n=m'/n'$ with both pairs coprime, then
$mn'=m'n$ implies $m\mid m'$ and $m'\mid m$, so $m=m'$ and $n=n'$.
This proves the final assertion without identifying differently decorated
or nonprimitive presentations.
\end{proof}

\begin{corollary}[Consecutive-shell intersection]
\label{cor:connes-consecutive}
If $b=a+1$ and both shells are in the allowed interval, write $R=R_a$.
Then the two shells have a common ordered primitive pair exactly when
\[
 R(R+4)\mid p+1.
\]
If this divisibility holds their common pairs are exactly $(p,1)$ and
$(1,p)$; otherwise the common set is empty.
For $(p,1)$ the exact denominator pair at shell $t=a,a+1$ is
\[
 \left(\frac{pt(p+1)}{4t-p},\frac{t(p+1)}{4t-p}\right),
\]
and the other primitive pair interchanges its entries.
\end{corollary}
\begin{proof}
Here $\gcd(a,a+1)=1$, so the embedded intersection is $H_p$ and the
retained integer divisors are $1,p$.  Its coprime pairs are
$(1,1),(1,p),(p,1)$.  The first has sum $2$, which is not divisible by
$R\ge3$.  The other two have sum $p+1$.
Also $\gcd(R,R+4)=\gcd(R,4)=1$ since $R$ is odd.  The preceding theorem
therefore imposes exactly $R(R+4)\mid p+1$.
For $(m,n)=(p,1)$, the formula for $k_t$ becomes
$t(p+1)/(4t-p)$; multiplying by $p$ or $1$ gives the displayed pair.
\end{proof}

\begin{proposition}[What the abstract-point isomorphism retains]
\label{prop:connes-abstract-loss}
For positive $N,M$, the unique positive ordered-group isomorphism
$H_N\to H_M$ is
\[
 \phi_{N,M}(x)=\frac NM x.
\]
On reciprocal divisors its exact integral-return subset and map are
\[
 u=N/e\longmapsto v=M/e
 \qquad(e\mid\gcd(N,M)).
\]
The inverse sends $M/e$ back to $N/e$; these fibres are singletons.
All finite embedded groups $H_N$ have the same abstract ordered-group
isomorphism class, so the fibre of that passage consists of every finite
positive budget $N$.
For the shells above, applying $\phi_{pa,pb}$ coordinatewise to the
reciprocal of a primitive pair yields the reciprocal integer labels
\[
 (m',n')=\left(\frac ba m,\frac ba n\right).
\]
This is an integral primitive pair only when $a=b$.  In particular, for
two distinct shells this isomorphism cannot itself give a successor on
the actual primitive pairs, although their embedded intersection has the
complete correspondence and positive witness return just proved.
\end{proposition}
\begin{proof}
The least positive element of $H_N$ is $1/N$, and an ordered-group
isomorphism must send it to the least positive element $1/M$ of $H_M$.
Additivity forces the given formula, which is indeed an isomorphism.
For $u\mid N$ write $e=N/u$.  Then
$\phi_{N,M}(1/u)=e/M$ is a reciprocal positive integer $1/v$ exactly
when $e\mid M$, in which case $v=M/e\mid M$.
Together with $e\mid N$ this gives precisely the stated retained subset
and its inverse.  Existence of the isomorphism for every $N,M$ proves
the assertion about the isomorphism-class fibre; as this passage is a
map on embedded objects, its kernel relation is all pairs of finite
positive budgets, not the kernel of an additive homomorphism.
For shell pairs $\phi_{pa,pb}(1/m)=a/(bm)$, which gives the displayed
integer label $m'=bm/a$ if integral, and likewise for $n$.
Write $a=dk$, $b=d\ell$, where $d=\gcd(a,b)$ and $\gcd(k,\ell)=1$.
Integrality of both labels forces $k\mid m,n$.  Their original primitivity
therefore forces $k=1$.  Now $(m',n')=(\ell m,\ell n)$ has greatest
common divisor $\ell$, so its primitivity forces $\ell=1$.
Consequently $a=b=d$.  Conversely if $a=b$ the map is the identity
and retains every primitive pair.  This proves the exact obstruction for
this specific isomorphism, alongside the surviving embedded correspondence.
\end{proof}

\subsection{Composing the finite point with the middle-channel shear}
\label{subsec:connes-middle-shear}
For an allowed shell $a$ retain the full divisor set
$\mathcal D_{a^2}$ and the two marks
\[
 E_a=\{u\mid a^2:R_a\mid4u+1\},\qquad
 M_a=\{u\mid a^2:R_a\mid u+a\}.
\]
The following maps have different domains and will all be kept in the
composition: the ordered-group isomorphism on $H_{a^2}$, its reciprocal
divisor return, and a divisor-to-primitive-pair map depending on the shell.

\begin{proposition}[Full middle-channel primitive return]
\label{prop:connes-middle-return}
There is a bijection
\[
 \mu_a:\mathcal D_{a^2}\longrightarrow
 \{(m,n):m\mid a,\ n\mid a,\ \gcd(m,n)=1\},
\]
\[
 \mu_a(u)=\left(\frac{u}{\gcd(u,a)},\frac{a}{\gcd(u,a)}\right).
\]
Its inverse is $(m,n)\mapsto am/n$ and its fibres are singletons.
It preserves the exact mark by
\[
 R_a\mid u+a\quad\Longleftrightarrow\quad R_a\mid m+n.
\]
For a marked divisor its denominator return is
\[
 \left(\frac{p(a+u)}{R_a},\frac{p(a+a^2/u)}{R_a}\right),
\]
and the ordered ratio of these entries is exactly $u/a=m/n$.
\end{proposition}
\begin{proof}
Put $g=\gcd(u,a)$.  Then $m=u/g$, $n=a/g$ are coprime.
For a prime $\ell$ let $e=v_\ell(a)$ and $j=v_\ell(u)$.
Since $u\mid a^2$, $0\le j\le2e$.  The valuations of $m,n$ are,
respectively, $j-\min(j,e)$ and $e-\min(j,e)$, both in $[0,e]$;
at least one is zero.  Thus $m,n\mid a$ and the entire exponent data
are retained by the formulas.  Conversely, if $m,n\mid a$ are coprime,
$u=am/n$ is integral and $a^2/u=an/m$ is integral.  Writing $a=gn$
gives $u=gm$ and $\gcd(u,a)=g\gcd(m,n)=g$, proving both inverse identities.
Also $u+a=g(m+n)$ and $g\mid a$, so $\gcd(g,R_a)=1$; this proves
the mark equivalence.  Proposition~\ref{prop:connes-pair-witness}
applies to this primitive pair.  Its $k$ is
$pa(m+n)/(mnR_a)$, and substitution of $u=am/n$ shows that
$km=p(a+u)/R_a$ and $kn=p(a+a^2/u)/R_a$.
Thus both are positive integers giving the required shell witness.
Finally dividing the displayed entries, or using $(km)/(kn)$, gives
$m/n=u/a$ exactly.
\end{proof}

\begin{theorem}[Exact marked domain of the arithmetic-site successor]
\label{thm:connes-marked-successor}
Let $b=a+1$ be another allowed shell and put $R=4a-p$.
The ordered-group isomorphism $H_{a^2}\to H_{(a+1)^2}$ has reciprocal
divisor return only at
\[
 u=a^2\longmapsto v=(a+1)^2.
\]
Neither endpoint is $E$-marked.  Its full domain of successful $E/M$
marked returns consists exactly of the middle-channel endpoints for
which
\[
 R(R+4)\mid p+4.
\]
For every such endpoint the maps compose as
\[
 \begin{array}{ccc}
 a^2&\longmapsto&(a+1)^2\\
 \mu_a\downarrow&&\downarrow\mu_{a+1}\\
 (a,1)&\xrightarrow{\ L\ }&(a+1,1),
 \end{array}
 \qquad
 L=\begin{pmatrix}1&1\\0&1\end{pmatrix},
\]
where the lower arrow acts on column vectors.
The exact positive integer denominator map is
\[
 \left(a,\frac{pa(a+1)}R,\frac{p(a+1)}R\right)
 \longmapsto
 \left(a+1,\frac{p(a+1)(a+2)}{R+4},\frac{p(a+2)}{R+4}\right).
\]
Both triples satisfy the Erd\H{o}s--Straus equation for the unchanged
prime $p$.  The inverse on this marked image retains $p$ and sends
$(a+1,R+4,(a+1)^2)$ back to $(a,R,a^2)$; its fibres are singletons.

Equivalently, all successful parameter triples $(p,R,a)$ have the form
\[
 \begin{split}
 &R\ge3,\quad R\equiv3\pmod4,\quad t\in\mathbb Z_{\ge0},\\
 &p=R(R+4)(4t+1)-4,\qquad p\text{ prime},\qquad p\equiv1\pmod{12},\\
 &a=\frac{R(R+5)}4-1+R(R+4)t.
 \end{split}
\]
For these parameters both $a$ and $a+1$ are automatically in the
prescribed shell interval.
\end{theorem}
\begin{proof}
The reciprocal divisor return calculation in
Proposition~\ref{prop:connes-abstract-loss}, applied to
$N=a^2$, $M=(a+1)^2$ with greatest common divisor $1$, gives exactly
the top-divisor map.
We prove the $E$ assertion without dropping the sign or the factor $4$.
Any positive odd integer $R\equiv3\pmod4$ has a prime divisor
$q\equiv3\pmod4$: otherwise every prime factor would be $1$ modulo
$4$, or every occurrence of a $3$-modulo-$4$ prime would have even
exponent, and the product would be $1$ modulo $4$.
If $R\mid4a^2+1$, then modulo this $q$ the nonzero residue $x=2a$
satisfies $x^2=-1$.  The residues $1,x,-1,-x$ are distinct and form a
subgroup of order $4$ in $(\mathbb Z/q\mathbb Z)^\times$.
Its cosets partition the $q-1$ nonzero residues into sets of four,
forcing $4\mid q-1$, a contradiction.  Thus $R\nmid4a^2+1$.
The same proof with $a+1,R+4$ excludes the next $E$ mark.
The middle mark at the first endpoint is
\[
 R\mid a^2+a\quad\Longleftrightarrow\quad R\mid a+1
 \quad\Longleftrightarrow\quad R\mid p+4.
\]
The first equivalence uses $\gcd(a,R)=1$; the second uses
$4(a+1)=p+R+4$ and $\gcd(4,R)=1$.
At the next endpoint the same identities give
$R+4\mid(a+1)^2+(a+1)$ exactly when $R+4\mid a+2$,
equivalently $R+4\mid p+4$.
Since $\gcd(R,R+4)=1$, both marks hold exactly when
$R(R+4)\mid p+4$.  This calculates the successful domain; no lifting
condition is left unevaluated.
Proposition~\ref{prop:connes-middle-return} gives
$\mu_a(a^2)=(a,1)$ and $\mu_{a+1}((a+1)^2)=(a+1,1)$.
Multiplication by $L$ sends $(a,1)^{\mathsf T}$ to
$(a+1,1)^{\mathsf T}$, so the displayed composition commutes.
The denominator formulas are the middle-channel return at these two
divisors, and that proposition proves their integrality, positivity,
and reciprocal identity.  Retaining $p,a,R$ gives the stated inverse.
This composition is not the coordinatewise map $\phi_{pa,pb}$ on
primitive reciprocal pairs: it changes the ratio from $a$ to $a+1$
through the two different shell-dependent maps $\mu_a,\mu_{a+1}$.
Consequently it does not contradict the no-return assertion for that
different map in Proposition~\ref{prop:connes-abstract-loss}.

For the parameter description put $k=(p+4)/(R(R+4))$.
It is a positive integer; both $p+4$ and $R(R+4)$ are $1$ modulo
$4$, so $k\equiv1\pmod4$ and uniquely $k=4t+1$ with $t\ge0$.
Then $p=R(R+4)(4t+1)-4$ and $a=(p+R)/4$ give the displayed formula.
Conversely those formulas imply $R=4a-p$ and the successful divisibility.
Set $h=(p-1)/12$.  Since $R\ge3$ and $R\equiv3\pmod4$,
$a=(p+R)/4\ge(p+3)/4=3h+1$ and is integral.
Also
\[
 2p-(R+7)=2R(R+4)(4t+1)-R-15
 \ge2R^2+7R-15>0\qquad(R\ge3).
\]
Thus $p+R+4\le3p-3$, which gives
$a+1\le(3p-3)/4=9h$.  This proves the required interval assertions.
\end{proof}

\paragraph{Exact successful example.}
For $R=7$ and $t=0$ the parameter formula gives $p=73$, $h=6$,
$a=20$, and $a+1=21$, both in $\{19,\ldots,54\}$.
The divisor map is $400\mapsto441$, the primitive map is
$(20,1)\mapsto(21,1)$, and the two triples are
\[
 (20,4380,219)\longmapsto(21,3066,146).
\]
Indeed $p+4=77=7\cdot11$,
$400+20=420=7\cdot60$, and $441+21=462=11\cdot42$.
Their exact reciprocal identities are
\[
 \frac1{20}+\frac1{4380}+\frac1{219}=\frac4{73},\qquad
 \frac1{21}+\frac1{3066}+\frac1{146}=\frac4{73},
\]
as follows by putting the first sum over denominator $4380$,
giving $(219+1+20)/4380=240/4380=4/73$, and the second over
$3066$, giving $(146+1+21)/3066=168/3066=4/73$.
